\subsection{\texorpdfstring{同态 \(\operatorname{Hom}_{\mathbb{C}}(\mathrm{E}_1, \mathrm{F}_1) \otimes_{\mathbb{C}} \operatorname{Hom}_{\mathbb{C}}(\mathrm{E}_2, \mathrm{F}_2) \to \operatorname{Hom}_{\mathbb{C}}(\mathrm{E}_1 \otimes_{\mathbb{C}} \mathrm{E}_2, \mathrm{F}_1 \otimes_{\mathbb{C}} \mathrm{F}_2)\)}{从 E1,F1 上的 Hom 与 E2,F2 上的 Hom 的 C 上张量积到相应张量积上的 Hom 的同态}}

设 C 为交换环，且 \(E_{1}\) , \(E_{2}\) , \(F_{1}\) , \(F_{2}\) 为四个 C-模；在 §3 第 5 号公式 (13) 中，我们定义了典范 C-模同态

\[
\lambda : \operatorname{Hom} (\mathbf {E} _ {1}, \mathbf {F} _ {1}) \otimes \operatorname{Hom} (\mathbf {E} _ {2}, \mathbf {F} _ {2}) \to \operatorname{Hom} (\mathbf {E} _ {1} \otimes \mathbf {E} _ {2}, \mathbf {F} _ {1} \otimes \mathbf {F} _ {2}). \tag {21}
\]

命题4. 当有序对 \((\mathbf{E}_1, \mathbf{E}_2)\)、\((\mathbf{E}_1, \mathbf{F}_1)\)、\((\mathbf{E}_2, \mathbf{F}_2)\) 之一由有限生成投射C-模组成时，典范同态 (21) 是双射。

显然，只需对有序对 \((\mathbf{E}_{1}, \mathbf{F}_{1})\) 与 \((\mathbf{E}_{1}, \mathbf{E}_{2})\) 证明即可。

我们先考虑有序对 \((\mathbf{E}_1, \mathbf{F}_1)\) 的情形；我们固定 \(\mathbf{E}_2, \mathbf{F}_1, \mathbf{F}_2\)，并对每个C-模记 \(\mathrm{T}(\mathbf{E}) = \mathrm{Hom}(\mathbf{E}, \mathbf{F}_1) \otimes_{\mathbb{C}} \mathrm{Hom}(\mathbf{E}_2, \mathbf{F}_2)\) 与 \(\mathrm{T}'(\mathbf{E}) = \mathrm{Hom}(\mathbf{E} \otimes \mathbf{E}_2, \mathbf{F}_1 \otimes \mathbf{F}_2)\)，且对每个C-同态 \(v: \mathbf{E} \to \mathbf{E}'\)，

\[
\mathrm{T} (v) = \operatorname{Hom} (v, \mathrm{l} _ {\mathrm{F} _ {1}}) \otimes \mathrm{l} _ {\operatorname{Hom} (\mathrm{E} _ {2}, \mathrm{F} _ {2})}
\]

和

\[
\mathrm{T} ^ {\prime} (v) = \operatorname{Hom} (v \otimes 1 _ {\mathbb {E} _ {2}}, 1 _ {\mathbf {F} _ {1} \otimes \mathbf {F} _ {2}}).
\]

于是引理4与引理5（第2号）（其中v换成了λ）成立，且可用完全类似的方法证明。

接下来我们固定 \(\mathbf{E}_2\) 与 \(\mathbf{F}_2\)，这次对每个C-模 \(\mathbf{F}\) 记 \(\mathrm{T}(\mathbf{F}) = \mathrm{Hom}(\mathbf{C},\mathbf{F})\otimes_{\mathbb{C}}\mathrm{Hom}(\mathbf{E}_2,\mathbf{F}_2)\) 与 \(\mathrm{T}'(\mathrm{F}) = \mathrm{Hom}(\mathbf{C}\otimes \mathbf{E}_2,\mathbf{F}\otimes \mathbf{F}_2)\)，且对每个C-同态 \(u:\mathrm{F}\to \mathrm{F}'\)，

\[
\mathrm{T} (u) = \operatorname{Hom} (\mathrm{l} _ {\mathbf {c}}, u) \otimes \mathrm{l} _ {\operatorname{Hom} (\mathbb {E} _ {2}, \mathrm{F} _ {2})}
\]

和

\[
\mathrm{T} ^ {\prime} (u) = \operatorname{Hom} (\mathrm{l} _ {\mathrm{c}} \otimes \mathrm{l} _ {\mathrm{E} _ {2}}, u \otimes \mathrm{l} _ {\mathrm{F} _ {2}}).
\]

这次可立即验证引理1与引理2（第2号）（其中 \(\lambda\) 处处代替 \(\nu\)）成立。

这样，我们先在 \(E_{1} = C\) 且 \(F_{1}\) 为投射且有限生成的情形证明该命题。第2号的论证（基于引理1与引理2）连同上述说明，把它归结为在 \(F_{1} = C\) 时证明该命题；于是 \(\operatorname{Hom}(E_{1}, F_{1})\)、\(E_{1} \otimes E_{2}\) 与 \(F_{1} \otimes F_{2}\) 分别等同于C、\(E_{2}\) 与 \(F_{2}\)（§3，第4号，命题4）；从而 (21) 式的两边都典范地等同于 \(\operatorname{Hom}(\mathbf{E}_{2},\mathbf{F}_{2})\)，且在这些等同之后，可验证 \(\lambda\) 成为恒等映射。

我们现在假设 \(F_{1}\) 是投射且有限生成的；第2号的论证（这次依赖引理4与引理5）把 \(E_{1}\) 为任意有限生成投射模时的证明归结为 \(E_{1} = C\) 的情形，即归结为已处理的第一种情形。

对于有序对 \((\mathbf{E}_{1}, \mathbf{E}_{2})\) ，过程类似，这次两次应用引理4和5；细节留给读者。

注意当 \(\mathbf{E}_1 = \mathbf{C}^{(I)}\)、\(\mathbf{E}_2 = \mathbf{C}^{(J)}\) 是自由模（不论是否有限生成）时，有 \(\operatorname{Hom}(\mathbf{E}_1, \mathbf{F}_1) = \mathbf{F}_1^{\mathrm{I}}\)、\(\operatorname{Hom}(\mathbf{E}_2, \mathbf{F}_2 = \mathbf{F}_2^{\mathrm{J}}\) 以及

\[
\operatorname{Hom} \left(\mathrm{E} _ {1} \otimes \mathrm{E} _ {2}, \mathrm{F} _ {1} \otimes \mathrm{F} _ {2}\right) = \left(\mathrm{F} _ {1} \otimes \mathrm{F} _ {2}\right) ^ {\mathrm{I} \times \mathrm{J}}
\]

（这在典范同构的意义下成立），且 (21) 此时就是§3第7号中典范同态 (22) 的一个特例。

当 \(E_{2} = C\) 时，把 \(\operatorname{Hom}(E_{2}, F_{2})\) 与 \(F_{2}\)、\(E_{1} \otimes E_{2}\) 与 \(E_{1}\) 等同后，典范同态 (21) 给出一个典范同态

\[
\operatorname{Hom} (\mathrm{E}, \mathrm{F}) \otimes \mathrm{G} \rightarrow \operatorname{Hom} (\mathrm{E}, \mathrm{F} \otimes \mathrm{G})\tag{22}
\]

对任意三个C-模E、F、G，它恰好是 \(\mathbf{A} = \mathbf{B} = \mathbf{C}\) 时第2号中的同态 (7)。

注意当 \(\mathbf{F} = \mathbf{C}\) 时，典范同态 (22) 再次给出交换环情形下的 (11)（第2号）。

现在假设 \(\mathbf{F}_1 = \mathbf{F}_2 = \mathbf{C}\)；由于 \(\mathbf{F}_1 \otimes \mathbf{F}_2\) 与 \(\mathbf{C}\) 等同，这次得到一个典范同态

\[
\mu : \mathbf {E} ^ {*} \otimes \mathbf {F} ^ {*} \rightarrow (\mathbf {E} \otimes \mathbf {F}) ^ {*}\tag{23}
\]

对于两个 C-模 E, F；对 \(x^{*} \in \mathbf{E}^{*}\)、\(y^{*} \in \mathbf{F}^{*}\)，\(x^{*} \otimes y^{*}\) 在典范同态 (23) 下的像是线性形式 \(u\)（\(\mathbf{E} \otimes \mathbf{F}\) 上的），满足

\[
u (x \otimes y) = \langle x, x ^ {*} \rangle \langle y, y ^ {*} \rangle .\tag{24}
\]

此外，若 \(E_{1}\) , \(E_{2}\) , \(F_{1}\) , \(F_{2}\) 是四个 C-模， \(f: E_{1} \rightarrow E_{2}\) , \(g: F_{1} \rightarrow F_{2}\) 是两个线性映射，则由 (24) 立即得出图表

\[
\begin{array}{c} \mathrm{E} _ {2} ^ {*} \otimes \mathrm{F} _ {2} ^ {*} \xrightarrow {\mu} (\mathrm{E} _ {2} \otimes \mathrm{F} _ {2}) ^ {*} \\ t _ {f} \otimes t _ {g} \Biggl \downarrow \qquad \qquad \qquad \qquad \qquad \qquad \qquad \qquad \qquad \qquad \qquad \qquad \qquad \qquad \qquad \qquad \qquad \qquad \qquad \qquad \qquad \qquad \qquad \qquad \qquad \qquad \qquad \qquad \qquad \qquad \qquad \qquad \qquad \qquad \\ \mathrm{E} _ {1} ^ {*} \otimes \mathrm{F} _ {1} ^ {*} \xrightarrow {\mu} (\mathrm{E} _ {1} \otimes \mathrm{F} _ {1}) ^ {*} \end{array}\tag{25}
\]

是交换的。

推论 1. 若模 E, F 中有一个是投射且有限生成的，则典范同态 (23) 是双射。

推论2. 设 \(\mathbf{E}_1, \mathbf{E}_2\) 是两个有限生成的投射C-模，\(u_1\) 是 \(\mathbf{E}_1\) 的一个自同态，\(u_2\) 是 \(\mathbf{E}_2\) 的一个自同态；则

\[
\operatorname{Tr} (u _ {1} \otimes u _ {2}) = \operatorname{Tr} (u _ {1}) \operatorname{Tr} (u _ {2}).\tag{26}
\]

由线性性，只需考虑 \(u_{1}\) 形如 \(x_{1} \mapsto \langle x_{1}, x_{1}^{*} \rangle y_{1}\) 且 \(u_{2}\) 形如 \(x_{2} \mapsto \langle x_{2}, x_{2}^{*} \rangle y_{2}\) 的情形；则 \(x_{1} \otimes x_{2}\) 在 \(u_{1} \otimes u_{2}\) 下的像按定义为

\[
\langle x _ {1}, x _ {1} ^ {*} \rangle \langle x _ {2}, x _ {2} ^ {*} \rangle (y _ {1} \otimes y _ {2}) = \langle x _ {1} \otimes x _ {2}, x _ {1} ^ {*} \otimes x _ {2} ^ {*} \rangle (y _ {1} \otimes y _ {2})
\]

其中 \(x_{1}^{*}\otimes x_{2}^{*}\) 在 \(\mu\) 下典范地等同于 \((\mathrm{E}_{1}\otimes\mathrm{E}_{2})^{*}\) 中的一个元素。由于 \(\langle y_{1}\otimes y_{2},x_{1}^{*}\otimes x_{2}^{*}\rangle=\langle y_{1},x_{1}^{*}\rangle\langle y_{2},x_{2}^{*}\rangle\)，公式 (26) 在此情形由 (17) 得出。

注记。若E、F、G是任意三个C-模，则容易验证图表

\[
\begin{array}{c} \mathrm{E} ^ {*} \otimes \mathrm{F} ^ {*} \otimes \mathrm{G} ^ {*} \xrightarrow {\mu \otimes 1} (\mathrm{E} \otimes \mathrm{F}) ^ {*} \otimes \mathrm{G} ^ {*} \\ 1 \otimes \mu \Biggl \downarrow \\ \mathrm{E} ^ {*} \otimes (\mathrm{F} \otimes \mathrm{G}) ^ {*} \xrightarrow [ \mu ]{} (\mathrm{E} \otimes \mathrm{F} \otimes \mathrm{G}) ^ {*} \end{array}\tag{27}
\]

是交换的，这可由公式(24)得出。

我们还注意到，在不假设C-模E、F有任何条件的情况下，存在典范同构

\begin{enumerate}
\def\labelenumi{(\arabic{enumi})}
\setcounter{enumi}{27}
\tightlist
\item
\end{enumerate}

\[
(\mathbf {E} \otimes \mathbf {F}) ^ {*} \rightarrow \operatorname{Hom} (\mathbf {E}, \mathbf {F} ^ {*})\tag{29}
\]

\[
(\mathbf {E} \otimes \mathbf {F}) ^ {*} \rightarrow \operatorname{Hom} (\mathbf {F}, \mathbf {E} ^ {*})
\]

它们即第1号中取 \(\mathbf{G} = \mathbf{C}\)、\(\mathbf{A} = \mathbf{B} = \mathbf{C}\) 时的同构 (6) 与 (5)。

这样，在 \(E \times F\) 上的双线性形式、从E到 \(F^{*}\) 的同态与从F到 \(E^{*}\) 的同态之间，就定义了一个典范的一一对应：若 u（相应地 v）是从E到 \(F^{*}\)（相应地，从F到 \(E^{*}\)）的同态，则相应的双线性形式由下式给出

\[
(x, y) \mapsto \langle y, u (x) \rangle \quad (\text { resp. } (x, y) \mapsto \langle x, v (y) \rangle).
\]

